% Compile with XeLaTeX from the repository root. This source references the Q3 output figures.
% Options for packages loaded elsewhere
\PassOptionsToPackage{unicode}{hyperref}
\PassOptionsToPackage{hyphens}{url}
%
\documentclass[UTF8,12pt,a4paper,fontset=windows]{ctexart}
\usepackage{amsmath,amssymb}
\usepackage{lmodern}
\usepackage{iftex}
\ifPDFTeX
  \usepackage[T1]{fontenc}
  \usepackage[utf8]{inputenc}
  \usepackage{textcomp} % provide euro and other symbols
\else % if luatex or xetex
  \usepackage{unicode-math}
  \defaultfontfeatures{Scale=MatchLowercase}
  \defaultfontfeatures[\rmfamily]{Ligatures=TeX,Scale=1}
  \setmainfont[]{SimSun}
\fi
% Use upquote if available, for straight quotes in verbatim environments
\IfFileExists{upquote.sty}{\usepackage{upquote}}{}
\IfFileExists{microtype.sty}{% use microtype if available
  \usepackage[]{microtype}
  \UseMicrotypeSet[protrusion]{basicmath} % disable protrusion for tt fonts
}{}
\makeatletter
\@ifundefined{KOMAClassName}{% if non-KOMA class
  \IfFileExists{parskip.sty}{%
    \usepackage{parskip}
  }{% else
    \setlength{\parindent}{0pt}
    \setlength{\parskip}{6pt plus 2pt minus 1pt}}
}{% if KOMA class
  \KOMAoptions{parskip=half}}
\makeatother
\usepackage{xcolor}
\IfFileExists{xurl.sty}{\usepackage{xurl}}{} % add URL line breaks if available
\IfFileExists{bookmark.sty}{\usepackage{bookmark}}{\usepackage{hyperref}}
\hypersetup{
  pdftitle={问题三：基于问题二前向机制的视觉认知动态宏观模型},
  pdflang={zh-CN},
  hidelinks,
  pdfcreator={LaTeX via pandoc}}
\urlstyle{same} % disable monospaced font for URLs
\usepackage[margin=2.5cm]{geometry}
\usepackage{longtable,booktabs,array}
\usepackage{calc} % for calculating minipage widths
% Correct order of tables after \paragraph or \subparagraph
\usepackage{etoolbox}
\makeatletter
\patchcmd\longtable{\par}{\if@noskipsec\mbox{}\fi\par}{}{}
\makeatother
% Allow footnotes in longtable head/foot
\IfFileExists{footnotehyper.sty}{\usepackage{footnotehyper}}{\usepackage{footnote}}
\makesavenoteenv{longtable}
\usepackage{graphicx}
\makeatletter
\def\maxwidth{\ifdim\Gin@nat@width>\linewidth\linewidth\else\Gin@nat@width\fi}
\def\maxheight{\ifdim\Gin@nat@height>\textheight\textheight\else\Gin@nat@height\fi}
\makeatother
% Scale images if necessary, so that they will not overflow the page
% margins by default, and it is still possible to overwrite the defaults
% using explicit options in \includegraphics[width, height, ...]{}
\setkeys{Gin}{width=\maxwidth,height=\maxheight,keepaspectratio}
% Set default figure placement to htbp
\makeatletter
\def\fps@figure{htbp}
\makeatother
\setlength{\emergencystretch}{3em} % prevent overfull lines
\providecommand{\tightlist}{%
  \setlength{\itemsep}{0pt}\setlength{\parskip}{0pt}}
\setcounter{secnumdepth}{3} % numbered through subsection level
\ifXeTeX
  % Load polyglossia as late as possible: uses bidi with RTL langages (e.g. Hebrew, Arabic)
  \usepackage{polyglossia}
  \setmainlanguage[]{}
\else
  \usepackage[main=]{babel}
% get rid of language-specific shorthands (see #6817):
\let\LanguageShortHands\languageshorthands
\def\languageshorthands#1{}
\fi
\ifLuaTeX
  \usepackage{selnolig}  % disable illegal ligatures
\fi

\title{问题三：基于问题二前向机制的视觉认知动态宏观模型}
\author{}
\date{}

\begin{document}
\maketitle

\hypertarget{ux6458ux8981}{%
\section*{摘要}\label{ux6458ux8981}}
\addcontentsline{toc}{section}{摘要}

针对题目要求的视觉认知建模与脑电验证，本文将问题二的视觉前向计算作为认知模型的输入，构造视觉状态 \(V\)、记忆相关状态 \(H\) 和控制状态 \(P\)，并用四份实测记录检验这些状态对 F3、Fz、F4 脑电波形的预测能力。数据以 VisCue 起点对齐，按统一质量规则从400个候选试次中保留355个。问题二模型生成的视觉群体活动驱动 \(V\)，提示门控、候选目标门控及假设匹配证据驱动 \(H\) 与 \(P\)；各状态参数预先固定，逐电极岭回归系数只在每折训练记录上估计。验证按完整记录留一，比较问题二视觉基线、加入 \(H\)、再加入 \(P\) 与交互项的预测误差。另将通道9每试次唯一的零到非零边沿作为候选应答时间参考 \(t_{\mathrm{act}}\)，逐试次把候选应答前信号截断于 \(t_{\mathrm{act}}-0.1\) s，并评价累计窗和末500 ms窗。

在固定晚期窗 \([0.8,2.8)\) s 中，加入 \(H\) 后仅因果滤波汇总误差微降0.03\%，无滤波和零相位滤波误差分别上升0.78\%和1.31\%；九个预处理×候选目标时刻汇总条件中仅两个误差下降。应答前窗的NRMSE约为1.02---1.06，加入记忆或控制后的改变量也很小。通道9标记中位于提示后2.2148 s，对应应答前端点中位数2.1148 s；355个纳入试次中354个端点早于2.2 s，全部早于2.4 s。因此，候选目标为2.2 s时354/355个试次的该窗不含目标出现后的脑电（另1个例外），候选目标为2.4 s时355个试次均不含目标后的脑电。另以独立的假设行为累积器模拟错误、迟答和漏答情景；这些结果只说明预设参数下能产生相应行为，不能证明真实EEG行为识别。本文完成了动态模型构造、跨记录留出检验和逐试次应答前探索性评价，但当前结果尚不支持记忆或控制状态具有稳定脑电预测增益，也不能据此确认具体脑区来源、任务正确率或实验意义上的迟答/漏答。

\textbf{关键词：} 视觉认知；脑电信号；动态宏观模型；记录留出验证；应答前窗口

\clearpage
\tableofcontents
\clearpage

\hypertarget{ux95eeux9898ux91cdux8ff0ux4e0eux5206ux6790}{%
\section{问题重述与分析}\label{ux95eeux9898ux91cdux8ff0ux4e0eux5206ux6790}}

题目要求在问题二脑电信号形成机制的基础上建立认知宏观模型，截取与视觉认知相关联的信号并用题目提供的脑电记录验证；题面指出信号终点可能在应答响应前约100 ms，并提醒应答可能不正确或未及时发生\cite{problem2026}。本研究的数据含四份连续记录、F3/Fz/F4脑电、VisCue视觉提示、通道9行动标记及时间戳。问题二提供的视觉处理和源到头皮前向投影为构造认知状态及电极观测关系提供了计算接口。

主要困难有三点。第一，三个额区电极不能唯一反演五个问题二源代理，因而新增状态不能直接解释成具体脑区活动。第二，通道9边沿集中在提示后约2.215 s，但数据没有独立的逐试次目标起始标记；真实目标时刻、任务映射和正式应答截止规则无法据此唯一确定。第三，模型参数与刺激场景若根据留出脑电反复选择，会使验证结果偏乐观。为此，本文将状态参数和目标场景作为预先设定的候选条件，逐电极观测系数只用训练记录估计，整份记录留作外层验证。

团队的《问题三建模设计》提出在问题二脑电形成机制基础上增加认知宏观状态并用实测脑电检验\cite{designQ3}；本文据当前代码落实其中的动态状态和跨记录预测验证。本文要检验的是：在同一问题二视觉前向基线上，加入记忆相关和控制状态能否改善未参与拟合记录的脑电预测。该检验反映候选状态的增量预测能力，不等同于认知机制的因果证明。

\hypertarget{ux5efaux6a21ux5047ux8bbe}{%
\subsection{建模假设}\label{ux5efaux6a21ux5047ux8bbe}}

\begin{enumerate}
\def\labelenumi{\arabic{enumi}.}
\tightlist
\item
  通道8 VisCue的零到非零边沿可重复地确定提示起点，仅用于提示锁定；它不替代独立目标起始时间。
\item
  目标类型、候选目标时刻、0.2 s持续时间和匹配证据是并列的机制情景，不是逐试次真实刺激或行为标签。
\item
  问题二默认参数、驱动尺度及V/H/P时间常数和增益在分析前固定；仅逐电极观测映射在每个训练折估计。
\item
  四份MAT文件作为外层留出记录，但文件与受试者的对应关系未确认，因此不把记录留出称为受试者独立验证。
\item
  V/H/P是功能性宏观状态，不能仅凭三个额区电极将其定位为海马、视觉皮层或前额叶信号。
\item
  通道9首次零到非零边沿仅作为候选应答时间参考；探索性窗口终点按该标记减100 ms计算，其与真实动作开始的同步关系尚未核实。
\end{enumerate}

\hypertarget{ux4e3bux8981ux7b26ux53f7}{%
\subsection{主要符号}\label{ux4e3bux8981ux7b26ux53f7}}

\begin{itemize}
\tightlist
\item
  \(t_i^c\)：第\(i\)试次VisCue提示起点；\(\tau\)：相对提示时间。
\item
  \(\mathbf y_i(\tau)\)：第\(i\)试次F3、Fz、F4实测脑电向量；\(\mathbf q_{Q2,k}\)：问题二五维源代理；\(G\)：五源代理到三电极的正向导联矩阵。
\item
  \(u_{Q2,k}\)：问题二早期视觉兴奋群体的空间均值；\(V_k,H_k,P_k\)：视觉、记忆相关和控制功能状态。
\item
  \(C_k,T_k\)：提示与候选目标门控；\(r\)：候选匹配证据；\(\rho_j,\tau_j\)：状态保持率与时间常数。
\item
  \(\alpha\)：岭回归惩罚系数；NRMSE：按留出实测标准差归一化的均方根误差。
\end{itemize}

\hypertarget{ux6570ux636eux4e8bux4ef6ux4e0eux9884ux5904ux7406}{%
\section{数据、事件与预处理}\label{ux6570ux636eux4e8bux4ef6ux4e0eux9884ux5904ux7406}}

\hypertarget{ux4e8bux4ef6ux4e0eux8bd5ux6b21}{%
\subsection{事件与试次}\label{ux4e8bux4ef6ux4e0eux8bd5ux6b21}}

四份原始记录均以256 Hz采样，每份记录含100个VisCue提示事件，共400个候选试次。以通道8 VisCue的零到非零边沿作为提示起点 \(t_i^c\)，从F3、Fz、F4连续信号构造提示相对区间 \([-0.2,3.0]\) s。处理后的试次统一重采样至128 Hz，并以提示前 \([-0.2,0)\) s均值作基线扣除。

质量控制剔除超出记录边界或与下一提示重叠的区间、含非有限值、原始幅度绝对值不小于999.5的区间，以及任一脑电通道平直的区间。原始单位未经独立校准，幅值门槛因此按数据原始单位描述，不换算为微伏。最终保留355个完整EEG试次，覆盖全部四份记录；固定窗与应答前窗评分均使用同一纳入集合。

通道8（VisCue）用于线索方向标签（−1/+1）；通道9（Action/TgtAct）记录应答方向编码（依数据标签解码为左右方向）。本次事件审计得到400个提示区间均有一个通道9边沿；提示与响应方向一致293次、不一致107次。这些计数描述的是线索---响应方向关系。由于Task-2线索到目标位置的正式映射未由实验记录核实，本文不把它们称为答对/答错次数；在缺少正式应答期限时，也不把区间内有无标记解释成及时率或漏答率。

\hypertarget{ux901aux90539ux6807ux8bb0ux8d77ux70b9ux5019ux9009ux5e94ux7b54ux65f6ux95f4ux53c2ux8003ux4e0eux5206ux6790ux7a97}{%
\subsection{通道9标记起点：候选应答时间参考与分析窗}\label{ux901aux90539ux6807ux8bb0ux8d77ux70b9ux5019ux9009ux5e94ux7b54ux65f6ux95f4ux53c2ux8003ux4e0eux5206ux6790ux7a97}}

对每个试次，若提示区间内通道9恰有一个零到非零边沿，则将该边沿记为候选应答时间参考 \(t_{\mathrm{act},i}\)。题目要求的候选应答前端点按

\[
t_{\mathrm{end},i}=t_{\mathrm{act},i}-0.100\text{ s}
\]

计算。该参考时刻用于逐试次信号截取和评分；通道9不作为状态输入、脑电预测变量或拟合目标。由于没有独立动作传感器或同步行为日志，通道9首次边沿是否与实际动作开始严格同步仍需核实，故不把它称作已确认的绝对应答时刻。

通道9边沿相对提示时间的中位数为2.2148 s，5\%---95\%分位数为2.2148---2.2188 s。故应答前端点中位数为2.1148 s，5\%---95\%分位数为2.1148---2.1188 s。逐试次审计还表明，355个有效端点中354个早于2.2 s，另1个端点为2.2281 s；全部端点早于2.4 s。于是，若候选目标时刻取2.2 s，多数试次的应答前窗在候选目标出现前已经结束，但有1个例外；若取2.4 s，则所有试次的端点均早于候选目标。2.0、2.2、2.4 s均作为敏感性候选值保留，不把任何一个当作已确认的真实目标时刻。

\begin{figure}
\centering
\includegraphics{src/C/q3/output/dynamic_heldout_validation/figures/analysis_event_windows.png}
\caption{事件标记、候选目标时刻与分析窗口}
\end{figure}

图1上部显示400个提示区间的通道9标记集中在约2.215 s；下部并列给出提示阶段、固定晚期评分窗、候选目标时刻和按通道9标记减100 ms得到的端点范围。端点中位数约2.115 s，早于2.2 s候选目标；逐试次审计中354/355个端点早于2.2 s，另1个例外，全部早于2.4 s。因而2.2和2.4 s候选目标情景的应答前窗一般不能包含目标出现后的加工。该图只显示时间对齐关系，不能据此确认候选目标是真实目标时刻，也不能证明通道9标记等同于动作起点；实际评分使用逐试次自己的 \(t_{\mathrm{act},i}-0.1\) s端点。

\hypertarget{ux9884ux5904ux7406}{%
\subsection{预处理}\label{ux9884ux5904ux7406}}

采用三种并行处理口径：不滤波；四阶Butterworth 0.5---30 Hz单向因果滤波；同阶同频带的离线零相位滤波。滤波先作用于连续EEG，再截取试次，以降低短epoch边界初始化影响。问题二基函数和状态特征采用相同的预处理算子、重采样及基线流程。零相位滤波仅用于离线敏感性比较；因果滤波保留了处理方向性，但仍可能存在相位延迟。

\begin{figure}
\centering
\includegraphics{src/C/q3/output/dynamic_heldout_validation/figures/observed_three_channel_stages.png}
\caption{F3、Fz、F4提示锁定实测脑电}
\end{figure}

图2显示三个电极在提示前后具有不同的组均值波形，后半段组间离散度也较大。每条细线对应一个“记录×提示方向”组，粗线和阴影分别为八组均值及其±1个组间标准差；它们不是逐试次置信带，也不是受试者总体推断。图中阴影的0.8—2.8 s仅标示预设的固定晚期分析窗，竖虚线为未经事件日志核实的候选目标时刻，故波形形态不能单独归因于目标加工或特定脑区。

\hypertarget{ux95eeux9898ux4e8cux524dux5411ux673aux5236ux63a5ux53e3}{%
\section{问题二前向机制接口}\label{ux95eeux9898ux4e8cux524dux5411ux673aux5236ux63a5ux53e3}}

问题二前向模型将视觉输入经视觉前端、LGN处理、皮层兴奋/抑制群体动力学及源代理映射，形成五个源代理的时序 \(\mathbf q_{Q2,k}\)，再通过导联矩阵 \(G\) 投影到三个电极：

\[
\mathbf y_{Q2,k}=G\mathbf q_{Q2,k},\qquad
\mathbf q_{Q2,k}\in\mathbb R^5,\quad
\mathbf y_{Q2,k}\in\mathbb R^3.
\]

这里沿用Wilson--Cowan兴奋/抑制群体相互作用的前向建模思想\cite{wilsonCowan}。三电极少于五个源代理，逆问题不唯一；本文只用 \(G\) 作正向投影，不求五源逆解。模型输入取问题二早期视觉兴奋群体的空间均值：

\[
u_{Q2,k}=\frac{1}{N_xN_y}\sum_{x=1}^{N_x}\sum_{y=1}^{N_y}
E_{\mathrm{early}}(x,y,k).
\]

本次调用问题二 \texttt{ModelParams} 的固定默认值 \((\tau_s,g_i,\tau_a)=(40\text{ ms},1.0,80\text{ ms})\)，视觉驱动尺度由固定镜像参考视觉刺激校准；问题三未使用四份实测EEG重新拟合这些问题二参数。实现还逐点校验源代理经 \(G\) 投影与问题二电极轨迹的一致性，最大绝对差为 \(5.95\times10^{-8}\)。该结果验证计算接口，不表示源代理可由三个电极唯一识别，也不证明前向模型具有生理真实性。

\hypertarget{ux89c6ux89c9ux8bb0ux5fc6ux63a7ux5236ux52a8ux6001ux6a21ux578b}{%
\section{视觉---记忆---控制动态模型}\label{ux89c6ux89c9ux8bb0ux5fc6ux63a7ux5236ux52a8ux6001ux6a21ux578b}}

\hypertarget{ux72b6ux6001ux542bux4e49ux548cux95e8ux63a7}{%
\subsection{状态含义和门控}\label{ux72b6ux6001ux542bux4e49ux548cux95e8ux63a7}}

定义 \(V_k\) 为受问题二视觉驱动的快速视觉状态，\(H_k\) 为提示写入并受候选匹配证据调制的记忆相关状态，\(P_k\) 为候选目标阶段由记忆强度及冲突情景驱动的控制状态。状态名称表示可检验的宏观功能假设，不分别等同于视觉皮层、海马或前额叶的实测活动。

设 \(C_k,T_k\in\{0,1\}\) 分别表示提示和候选目标门控，\(r\in[-1,1]\) 表示情景设定的匹配证据，正值为匹配、负值为不匹配。本次只取 \(r=+1\) 与 \(r=-1\)，它不是由通道8、通道9或真实任务结果生成的正确性标签。数值步长为 \(\Delta t=0.004\) s，状态按指数保持率

\[
\rho_j=\exp(-\Delta t/\tau_j),\qquad j\in\{V,H,P\}
\]

递推：

\[
\begin{aligned}
V_k&=\rho_VV_{k-1}+(1-\rho_V)u_{Q2,k},\\
H_k&=\rho_HH_{k-1}+(1-\rho_H)\left(g_sV_kC_k+g_mV_kT_kr\right),\\
P_k&=\rho_PP_{k-1}+(1-\rho_P)\left(g_{HP}|H_{k-1}|T_k+g_\delta T_k\frac{1-r}{2}\right).
\end{aligned}
\]

初值为 \(V_0=u_{Q2,0},H_0=P_0=0\)。其中 \(V_k\) 跟随问题二视觉输入；提示门控使 \(V_k\) 写入 \(H_k\)；目标门控及匹配情景进一步调节 \(H_k\)，并使记忆强度与不匹配项影响 \(P_k\)。按代码更新顺序，\(P_k\) 使用前一时刻 \(H_{k-1}\)。

\textbf{表1 动态状态的预设参数}

\begin{longtable}[]{@{}lrl@{}}
\toprule
参数 & 数值 & 模型作用 \\
\midrule
\endhead
\(\tau_V\) & 0.060 s & 视觉状态时间常数 \\
\(\tau_H\) & 1.000 s & 记忆相关状态时间常数 \\
\(\tau_P\) & 0.250 s & 控制状态时间常数 \\
\(g_s\) & 0.80 & 提示写入增益 \\
\(g_m\) & 0.80 & 目标匹配调制增益 \\
\(g_{HP}\) & 0.25 & 记忆至控制的耦合增益 \\
\(g_\delta\) & 0.80 & 不匹配/冲突情景增益 \\
\(g_{\mathrm{top}}\) & 0.10 & 构造性传感器示例中的控制调制增益 \\
\bottomrule
\end{longtable}

上述参数为预设机制值，没有用四份真实脑电或留出误差调参。构造性传感器方程为

\[
\mathbf y_{\mathrm{scenario},k}=(1+g_{\mathrm{top}}P_k)\mathbf y_{Q2,k}+\boldsymbol\ell_HH_k+\boldsymbol\ell_PP_k,
\]

其中加载向量取问题二固定的观测模式，用于生成机制示例，不是从三电极反演出的解剖导联。主要留出验证不直接固定这些加载，而是对Q2、H、P及Q2×P特征另行估计训练折内的电极映射。候选场景包含三种目标刺激（点阵、向内运动、向外运动），目标时刻为无目标、2.0、2.2、2.4 s，非空目标持续时间设为0.2 s；匹配证据取 ±1。目标刺激、时刻和匹配值均是候选情景，不是已核实的逐试次实验条件。

\begin{figure}
\centering
\includegraphics{src/C/q3/output/dynamic_heldout_validation/figures/dynamic_states_sensor_contributions.png}
\caption{一个候选情景下的V/H/P状态与电极贡献示例}
\end{figure}

图3上排给出一个点阵、2.2 s目标、不匹配分支下的视觉状态 \(V\)、记忆相关状态 \(H\) 和控制状态 \(P\)；下排显示训练折估计的状态/视觉输入观测成分与留出实测脑电的示例。该例中状态贡献曲线幅度明显小于部分留出波形变化，提示这些预设状态及映射不足以重构该代表性记录。它是单个留出记录和提示侧的机制示例，不是跨记录效果估计；成分加载是传感器空间系数，不能解释为海马或前额叶的解剖定位。

\hypertarget{ux7535ux6781ux89c2ux6d4bux6a21ux578bux4e0eux9a8cux8bc1ux65b9ux6cd5}{%
\section{电极观测模型与验证方法}\label{ux7535ux6781ux89c2ux6d4bux6a21ux578bux4e0eux9a8cux8bc1ux65b9ux6cd5}}

\hypertarget{ux9010ux7535ux6781ux89c2ux6d4bux65b9ux7a0b}{%
\subsection{逐电极观测方程}\label{ux9010ux7535ux6781ux89c2ux6d4bux65b9ux7a0b}}

对每个电极 \(c\in\{F3,Fz,F4\}\)，以预处理后的问题二电极轨迹 \(q_{c,k}\)、状态 \(H_k,P_k\) 及交互项 \(q_{c,k}P_k\) 构造递增特征。三个动态模型为

\[
\begin{aligned}
\mathbf x_{c,k}^{(0)}&=[q_{c,k}],\\
\mathbf x_{c,k}^{(H)}&=[q_{c,k},H_k],\\
\mathbf x_{c,k}^{(H,P)}&=[q_{c,k},H_k,P_k,q_{c,k}P_k].
\end{aligned}
\]

分别对应问题二视觉基线 \(Q2\)、\(Q2+H\) 和 \(Q2+H+P\)。训练折内对每个特征按均值和标准差标准化，目标电位按训练均值中心化，再用岭回归估计：

\[
\hat\beta_c=\arg\min_{\beta}
\left\|\mathbf y_{c,\mathrm{train}}-\bar y_{c,\mathrm{train}}\mathbf 1-
\mathbf Z_{c,\mathrm{train}}\beta\right\|_2^2+
\alpha\|\beta\|_2^2,\qquad \alpha=1,
\]

\[
\hat y_{c,k}=\bar y_{c,\mathrm{train}}+
\mathbf z_{c,k}^{\mathsf T}\hat\beta_c.
\]

岭回归用于在相关特征下约束系数波动\cite{hoerlKennard}。每折的特征均值、标准差、目标均值及回归系数只由训练记录计算，再原样用于留出记录。另设训练集同提示侧平均波形作为模板基线。

\hypertarget{ux7559ux51faux65b9ux6848ux4e0eux8befux5deeux6307ux6807}{%
\subsection{留出方案与误差指标}\label{ux7559ux51faux65b9ux6848ux4e0eux8befux5deeux6307ux6807}}

外层采用四折整记录留一：每次完整留出一份MAT文件，其余三份用于训练。训练单位为通过质量筛选的试次按``记录×提示侧''计算的均值波形；同一记录不拆分至训练和测试。每个测试组将三电极与窗内时间点展平，计算

\[
\mathrm{RMSE}_g=\sqrt{\frac{1}{n_g}\sum_{j=1}^{n_g}(y_{g,j}-\hat y_{g,j})^2},\qquad
\mathrm{NRMSE}_g=\frac{\mathrm{RMSE}_g}{\mathrm{SD}(y_g)},
\]

以及记忆增量指标

\[
\Delta_H=\mathrm{NRMSE}_{Q2}-\mathrm{NRMSE}_{Q2+H}.
\]

\(\Delta_H>0\) 表示加入 \(H\) 后误差下降。NRMSE按该留出组实测信号标准差归一化，越小表示相对误差越低。主结果使用固定晚期窗 \([0.8,2.8)\) s；早期窗 \([0,0.8)\) s和候选目标后局部窗用于补充诊断。所有候选刺激、候选目标时刻和预处理口径并列报告，不按测试误差挑选所谓真实目标时刻。

\hypertarget{ux9010ux8bd5ux6b21ux5e94ux7b54ux524dux8bc4ux4ef7}{%
\subsection{逐试次应答前评价}\label{ux9010ux8bd5ux6b21ux5e94ux7b54ux524dux8bc4ux4ef7}}

应答前分析只对通过EEG质控且在提示区间内有唯一通道9边沿的测试试次执行。累计窗为 \([0,t_{\mathrm{act},i}-0.1)\) s，末500 ms窗为 \([\max(0,t_{\mathrm{act},i}-0.6),t_{\mathrm{act},i}-0.1)\) s。先按每个试次自身的端点裁切实测EEG，再把同一留出``记录×提示侧''组内的试次样本拼接计算RMSE、NRMSE和相关系数；组级预测波形在各试次对应区间重复参与比较。组内计分试次数为41---49，中位数44.5。CSV所列窗起止时间只是逐试次边界的中位数摘要，评分本身没有改用公共端点。

这里需要区分模型拟合与评价：训练折的电极映射仍基于完整提示锁定组均值波形拟合；应答前端点用于留出集的逐试次截断评分，不用于重拟合一个只看应答前数据的模型。因此，应答前结果回答的是``全窗拟合出的候选模型在逐试次应答前数据上的预测误差如何''，不能表述成``模型参数完全由应答前脑电估计''。

\hypertarget{ux7ed3ux679cux4e0eux5206ux6790}{%
\section{结果与分析}\label{ux7ed3ux679cux4e0eux5206ux6790}}

\hypertarget{ux56faux5b9aux665aux671fux7a97ux53e3ux7684ux6574ux8bb0ux5f55ux7559ux51faux7ed3ux679c}{%
\subsection{固定晚期窗口的整记录留出结果}\label{ux56faux5b9aux665aux671fux7a97ux53e3ux7684ux6574ux8bb0ux5f55ux7559ux51faux7ed3ux679c}}

\textbf{表2 固定晚期窗 \([0.8,2.8)\) s的留出NRMSE}（均值，记录间标准差；每种处理4份留出记录）

\begin{longtable}[]{@{}lrrrr@{}}
\toprule
预处理 & 训练均值模板 & Q2 & Q2+H & Q2+H+P \\
\midrule
\endhead
无滤波 & 2.4744 (0.9864) & 2.3328 (0.8938) & 2.3511 (0.9131) & 2.3699 (0.9331) \\
因果滤波 & 2.1489 (1.6467) & 1.6479 (0.7586) & 1.6474 (0.7606) & 1.6834 (0.8317) \\
零相位滤波 & 2.0805 (1.4277) & 1.6144 (0.6406) & 1.6356 (0.6674) & 1.6183 (0.6549) \\
\bottomrule
\end{longtable}

因果滤波下 \(Q2+H\) 相对 \(Q2\) 的NRMSE从1.6479降至1.6474，相对改善约0.03\%；无滤波和零相位滤波下分别恶化约0.78\%和1.31\%。加入 \(P\) 后三个预处理口径均未优于对应的 \(Q2\) 视觉基线。九个``预处理×非空目标时刻''汇总单元中，ΔH仅有两个为正，二者均为因果滤波情景；其余七个单元误差未下降。由此只能报告个别情景出现很小的误差变化，不能认为记忆或控制动态获得稳定支持。

三种处理下 \(Q2\) 视觉基线均优于训练均值模板，表明问题二前向轨迹提供了可用于构造预测的结构。但表2所有NRMSE均大于1，留出预测误差仍达到或超过实测信号自身标准差。图4展示模型误差比较；图5给出一组留出波形作形态参照，不能替代整体统计结果。

\begin{figure}
\centering
\includegraphics{src/C/q3/output/dynamic_heldout_validation/figures/heldout_model_comparison.png}
\caption{不同模型的整记录留出误差比较}
\end{figure}

图4按预处理方式分面，比较训练均值模板、问题二视觉基线及增加记忆/控制状态后的早期提示窗与固定晚期窗平均NRMSE；柱高越低表示留出误差越小。问题二视觉基线在三种预处理下均优于训练均值模板，但加入记忆或控制后没有一致下降，且表2中的晚期窗NRMSE仍均大于1。图中是跨候选目标情景与留出记录的汇总均值，不能代替逐记录误差或证明新增状态稳定提升预测。

\begin{figure}
\centering
\includegraphics{src/C/q3/output/dynamic_heldout_validation/figures/heldout_predictions_vs_measured.png}
\caption{代表性留出记录的预测与实测波形}
\end{figure}

图5展示因果滤波、点阵刺激、2.2 s候选目标这一代表性情景：颜色区分四份留出记录，线型区分实测波形、问题二视觉预测和加入记忆后的预测。部分电极和记录的实测峰谷幅度、时序与两种预测差异明显，说明整体NRMSE所揭示的拟合不足也体现在代表性波形上。该图对提示侧及匹配分支取均值，只用于观察形态，不代表所有候选情景，也不能据单一情景挑选目标时刻。

\textbf{表3 记忆状态的固定晚期窗目标时刻敏感性}（ΔH；正值表示误差下降）

\begin{longtable}[]{@{}lrrr@{}}
\toprule
预处理 & 2.0 s & 2.2 s & 2.4 s \\
\midrule
\endhead
因果滤波 & −0.0084 & +0.0048 & +0.0052 \\
无滤波 & −0.0251 & −0.0208 & −0.0089 \\
零相位滤波 & −0.0229 & −0.0169 & −0.0237 \\
\bottomrule
\end{longtable}

九个汇总条件中仅因果滤波的2.2 s与2.4 s情景为正，且幅度有限。候选目标时刻未经事件日志确认，因此这一敏感性结果不应用来倒推出目标真实出现时间。

\begin{figure}
\centering
\includegraphics{src/C/q3/output/dynamic_heldout_validation/figures/memory_target_time_sensitivity.png}
\caption{不同候选目标时刻下的记忆态误差变化}
\end{figure}

图6的纵轴为 \(\Delta_H=\mathrm{NRMSE}_{Q2}-\mathrm{NRMSE}_{Q2+H}\)：正值表示加入记忆状态后误差下降，虚线为零基准；点为四份留出记录的均值，误差线为记录间标准差。只有因果滤波下2.2 s和2.4 s候选时刻的均值略为正，其余多数条件为负，且部分误差线跨越零。改善方向随滤波和候选时刻变化，支持“增益不稳定”的判断；候选时刻尚未核实，记录数也有限，故该图不构成对真实目标加工时刻或总体效应的推断。

\hypertarget{ux5e94ux7b54ux524dux7a97ux53e3ux7684ux9010ux8bd5ux6b21ux7559ux51faux7ed3ux679c}{%
\subsection{应答前窗口的逐试次留出结果}\label{ux5e94ux7b54ux524dux7a97ux53e3ux7684ux9010ux8bd5ux6b21ux7559ux51faux7ed3ux679c}}

\textbf{表4 逐试次应答前窗口NRMSE}（均值，条件行间标准差；每个单元汇总152个``19个候选情景×4份留出记录×2种提示侧''评分行）

\begin{longtable}[]{@{}llrr@{}}
\toprule
预处理 & 模型 & 累计窗 NRMSE & 末500 ms窗 NRMSE \\
\midrule
\endhead
因果滤波 & 训练均值模板 & 1.0409 (0.0652) & 1.0291 (0.0313) \\
因果滤波 & Q2 & 1.0264 (0.0419) & 1.0591 (0.0936) \\
因果滤波 & Q2+H & 1.0266 (0.0420) & 1.0583 (0.0922) \\
因果滤波 & Q2+H+P & 1.0283 (0.0444) & 1.0565 (0.0898) \\
无滤波 & 训练均值模板 & 1.0596 (0.0537) & 1.0606 (0.0598) \\
无滤波 & Q2 & 1.0557 (0.0534) & 1.0505 (0.0526) \\
无滤波 & Q2+H & 1.0537 (0.0513) & 1.0494 (0.0518) \\
无滤波 & Q2+H+P & 1.0523 (0.0501) & 1.0492 (0.0517) \\
零相位滤波 & 训练均值模板 & 1.0491 (0.0922) & 1.0309 (0.0249) \\
零相位滤波 & Q2 & 1.0212 (0.0229) & 1.0306 (0.0262) \\
零相位滤波 & Q2+H & 1.0206 (0.0233) & 1.0284 (0.0240) \\
零相位滤波 & Q2+H+P & 1.0239 (0.0293) & 1.0251 (0.0218) \\
\bottomrule
\end{longtable}

表4各值约为1.02---1.06，动态模型与视觉基线的差异很小。部分条件略有下降、部分略有上升，没有显示记忆或控制状态在不同预处理与窗口下稳定改善应答前预测。152行包含重复的目标候选情景，不是152个独立受试者或独立实验；其行间标准差仅用于描述情景与留出组的离散程度。

\hypertarget{ux76eeux6807ux65f6ux5e8fux5bf9ux8ba4ux77e5ux89e3ux91caux7684ux9650ux5236}{%
\subsection{目标时序对认知解释的限制}\label{ux76eeux6807ux65f6ux5e8fux5bf9ux8ba4ux77e5ux89e3ux91caux7684ux9650ux5236}}

应答前评价按逐试次端点执行，解决了``用组内中位数代替每个试次截断''的问题；但它并不能弥补真实目标时刻缺失。以候选目标2.2 s为例，354/355个端点早于该时刻，因此这些试次的应答前数据不包含该候选目标出现后的加工；一个端点晚于2.2 s，是例外。候选目标2.4 s晚于所有纳入试次的端点。因此，表4可用于评价候选模型在应答前数据上的波形误差，但不能用2.2或2.4 s分支验证目标出现后的线索---目标匹配、冲突或控制过程。2.0 s分支虽早于多数应答前端点约0.115 s，仍只是未经事件记录确认的候选条件。

\hypertarget{ux6a21ux578bux8bc4ux4ef7ux4e0eux9002ux7528ux8fb9ux754c}{%
\subsection{假设性行为情景仿真（不代表真实数据分类）}

为检验题目提及的错误、迟答与未应答行为是否能由一个明确的决策规则产生，本文另设独立随机证据累积器。它不接收EEG、V/H/P状态或真实行为标签，不参与脑电回归及留出评分。令 \(X_k\) 为时刻 \(k\) 的带符号决策证据，离散更新为

\[
X_{k+1}=X_k+\mu\Delta t+\sigma\sqrt{\Delta t}\,\varepsilon_k,
\qquad \varepsilon_k\sim\mathcal N(0,1).
\]

初值为 \(X_0=0\)。当 \(X_k\) 首次达到 \(+\theta\) 或 \(-\theta\) 时分别输出右选或左选；首次越界时刻作为模拟决策时间。表5采用 \(\Delta t=0.01\) s、\(\theta=1\)、截止时间 \(T_d=3\) s、观察期 \(T_{obs}=8\) s、噪声尺度 \(\sigma=0.05\)，每种情景模拟1000次，使用PCG64随机数生成器，种子为20260926，三种情景按错误、迟答、漏答顺序共用同一随机流。正确方向设为正向（右选）。错误情景取 \(\mu=-0.65\)；迟答情景取 \(\mu=+0.20\)；漏答情景取 \(\mu=0\)。阈值、漂移、噪声、截止时间、观察期、时间步长及模拟次数均为建模假设，而非实验参数估计。

\textbf{表5 假设参数下的行为情景仿真结果}（每种情景1000次）

\begin{center}
\small
\begin{tabular}{@{}lrrrrl@{}}
\toprule
情景 & 漂移 \(\mu\) & 错向选择 & 超过3 s截止 & 观察期内未越界 & 越界时间中位数（5\%--95\%） \\
\midrule
错误 & $-0.65$ & 1000 & 0 & 0 & 1.54 s（1.40--1.71） \\
迟答 & $+0.20$ & 0 & 1000 & 0 & 4.99 s（4.17--5.91） \\
漏答 & $0.00$ & 0 & -- & 1000 & 未越界 \\
\bottomrule
\end{tabular}
\end{center}

在这组参数下，程序生成的三类结果符合预先设定的情景定义：错误情景到达与正确方向相反的边界；迟答情景在截止时间后才到达正确方向边界；漏答情景在观察期结束前没有到达任何边界。该结果仅说明这个独立累积器在所设参数下可以产生三类行为。由于参数未由真实记录估计，仿真不能证明V/H/P脑电模型识别了真实数据中的错误、迟答或漏答，也不能据此估计真实发生率。

\section{模型评价与适用边界}\label{ux6a21ux578bux8bc4ux4ef7ux4e0eux9002ux7528ux8fb9ux754c}}

本文完成的问题三计算链包括：问题二视觉前向机制输入、V/H/P状态递推、逐电极观测映射、按整份记录留一的固定晚期窗预测比较、参数和场景敏感性分析，以及使用通道9标记起点作为候选应答时间参考进行逐试次截断的探索性评价。主要系数的标准化和岭回归只用训练记录拟合，问题二参数及V/H/P状态参数均为预设值，未用四份真实EEG调优，避免把同一留出记录用于前端参数拟合和Q3测试。

本文0.8—2.8秒主评分窗口为固定晚期分析窗，不保证全部位于实际应答之前。应答前信号依通道9标记起点逐试次截断；组内中位数仅用于汇总端点，不替代逐试次边界。该分析仅作探索性结果，通道9边沿是否等于真实动作起点尚未核实，因而不保证真正满足每个试次的应答前100毫秒要求。按现有端点审计，候选目标时刻为2.2秒时354/355个试次的分析窗不含目标后的加工（1个例外）；候选目标为2.4秒时355个试次均不含目标后的加工。

当前脑电结果只支持``模型已构造并完成留出检验''。它们不支持``模型较好解释脑电''：NRMSE大于1，加入H的误差改善不稳定，加入P无稳定增益；应答前模型参数也不是仅由逐试次截断数据拟合。四份文件是否对应四名独立受试者未确认，故记录留出不能等同于受试者外部验证。

此外，目标类型和时刻、Task-2线索到目标的映射、逐试次真实目标起始时间、正式应答截止规则及通道9边沿与实际动作起点的同步关系均未由实验资料核实。因此本文不报告真实数据的任务正确率、实验迟答率、漏答率或传统目标到动作反应时。本文另设独立假设性证据累积器模拟三种情景，但它未由真实EEG或行为标签拟合，不能证明模型识别了真实数据中的错误、迟答或漏答。通道8与通道9的方向一致/不一致只是可复算的编码关系，不等同于任务正确/错误。三个头皮电极也不足以唯一分解五个源代理，H/P加载不能证明海马或前额叶定位，更不构成精神疾病诊断效能证据。

后续若取得逐试次目标起始、刺激与任务映射、正式反应期限和同步动作日志，应先锁定事件定义，再在每个训练折中仅使用逐试次应答前数据估计观测模型，并将完整记录作为测试集。只有当记忆/控制状态在未参与拟合的记录与预先指定分析窗中稳定降低误差，才有依据讨论其增量解释能力。

\hypertarget{ux7ed3ux8bba}{%
\section{结论}\label{ux7ed3ux8bba}}

\begin{enumerate}
\def\labelenumi{\arabic{enumi}.}
\tightlist
\item
  基于问题二前向机制构造了视觉状态 \(V\)、记忆相关状态 \(H\) 和控制状态 \(P\)，并通过训练折内逐电极岭回归映射到F3、Fz、F4。
\item
  四份记录整记录留一结果显示：固定晚期窗中加入 \(H\) 仅在因果滤波汇总下产生约0.03\%的微小改善；其余预处理和多数目标时刻情景未改善，加入 \(P\) 未显示稳定增益。所有晚期窗NRMSE均大于1。
\item
  已按通道9候选应答时间参考逐试次截断 \(t_{\mathrm{act}}-100\) ms，并计算累计窗和末500 ms窗留出误差；各模型NRMSE约为1.02---1.06，状态增量仍不稳定。该应答前结果只是受事件时序限制的探索性分析。
\item
  独立假设累积器在设定参数下分别产生1000次错向选择、1000次截止后响应和1000次观察期未越界；这只说明模拟规则可产生这些行为，不能作为真实数据的错误、迟答或漏答识别结果。
\item
  当前可得结论是问题三模型流程已运行并完成初步跨记录检验；尚无充分证据说明\(H/P\)状态稳定解释脑电，也不能据现有数据给出任务正确率、真实迟答/漏答率或脑区定位结论。
\end{enumerate}

\hypertarget{ux53c2ux8003ux6587ux732eux4e0eux590dux73b0ux5165ux53e3}{%
\section{参考文献与复现入口}\label{ux53c2ux8003ux6587ux732eux4e0eux590dux73b0ux5165ux53e3}}

脑电验证主实现入口为 \texttt{src/C/q3/dynamic\_validation.py}，状态模型位于 \texttt{src/C/q3/dynamic\_cognitive\_model.py}。行为情景仿真按本文给出的累积方程、PCG64随机数生成器、参数、种子20260926及每种情景1000次复现；三种情景顺次共用同一随机流，参数均为假设值。

\begin{thebibliography}{99}
\bibitem{problem2026} 第二十三届中国研究生数学建模竞赛组委会. 服务于脑机接口与精神性疾病诊断的脑电图计算模型：C题题目[Z]. 2026.（题目及通道语义依据项目提供的题目文件。）
\bibitem{designQ3} 项目组. 问题三建模设计[内部建模文档]. 2026.（用于说明模型设计来源。）
\bibitem{wilsonCowan} Wilson H R, Cowan J D. Excitatory and inhibitory interactions in localized populations of model neurons[J]. \textit{Biophysical Journal}, 1972, 12(1): 1--24. \href{https://pubmed.ncbi.nlm.nih.gov/4332108/}{PubMed记录}.
\bibitem{hoerlKennard} Hoerl A E, Kennard R W. Ridge regression: Biased estimation for nonorthogonal problems[J]. \textit{Technometrics}, 1970, 12(1): 55--67. \href{https://www.tandfonline.com/doi/abs/10.1080/00401706.1970.10488634}{Taylor \& Francis记录}.
\end{thebibliography}

\end{document}
